\documentclass[11pt,letterpaper]{article}
\usepackage[T1]{fontenc}
\usepackage[utf8]{inputenc}
\usepackage{lmodern}
\usepackage{amsmath,amssymb,amsthm,mathtools,mathrsfs}
\usepackage[margin=1in,headheight=15pt]{geometry}
\usepackage{microtype}
\usepackage{xcolor,booktabs,array,longtable}
\usepackage{enumitem}
\usepackage{fancyhdr}
\usepackage[most]{tcolorbox}
\usepackage{listings}
\usepackage{hyperref}
\usepackage[nameinlink,capitalize,noabbrev]{cleveref}
\definecolor{ink}{HTML}{183247}
\definecolor{accent}{HTML}{206F78}
\definecolor{soft}{HTML}{EFF5F6}
\hypersetup{colorlinks=true,linkcolor=accent,citecolor=accent,urlcolor=accent,
 pdftitle={Gonshor's Product-Birthday Conjecture: Ordinal-Support Series},
 pdfauthor={Research draft prepared with ChatGPT}}
\pagestyle{fancy}
\fancyhf{}
\fancyhead[L]{\small\color{ink}Surreal product birthdays}
\fancyhead[R]{\small\color{ink}Research draft / 20 September 2026}
\fancyfoot[C]{\thepage}
\renewcommand{\headrulewidth}{0.3pt}
\setlength{\parskip}{0.35em}
\setlength{\parindent}{1.25em}
\setcounter{tocdepth}{1}
\allowdisplaybreaks[2]
\newtheorem{theorem}{Theorem}[section]
\newtheorem{proposition}[theorem]{Proposition}
\newtheorem{lemma}[theorem]{Lemma}
\newtheorem{corollary}[theorem]{Corollary}
\theoremstyle{definition}
\newtheorem{definition}[theorem]{Definition}
\newtheorem{example}[theorem]{Example}
\theoremstyle{remark}
\newtheorem{remark}[theorem]{Remark}
\crefname{lemma}{Lemma}{Lemmas}
\crefname{proposition}{Proposition}{Propositions}
\crefname{corollary}{Corollary}{Corollaries}
\crefname{theorem}{Theorem}{Theorems}
\newcommand{\No}{\mathbf{No}}
\newcommand{\On}{\mathbf{On}}
\newcommand{\R}{\mathbb{R}}
\newcommand{\Z}{\mathbb{Z}}
\newcommand{\N}{\mathbb{N}}
\newcommand{\D}{\mathbb{D}}
\newcommand{\A}{\mathscr{A}}
\newcommand{\bd}{\ell}
\newcommand{\ns}{\mathbin{\oplus}}
\newcommand{\np}{\mathbin{\otimes}}
\newcommand{\ep}{\varepsilon}
\DeclareMathOperator{\otp}{otp}
\DeclareMathOperator{\lc}{lc}
\DeclareMathOperator{\supp}{supp}
\DeclareMathOperator{\wt}{wt}
\newcommand{\cof}{\mathrm{cf}}
\newtcolorbox{resultbox}[1][]{colback=soft,colframe=accent,boxrule=0.6pt,
 arc=1mm,left=3mm,right=3mm,top=2mm,bottom=2mm,#1}
\lstset{basicstyle=\ttfamily\small,breaklines=true,columns=fullflexible,
 keepspaces=true,frame=single,rulecolor=\color{accent},showstringspaces=false}

\begin{document}
\hypersetup{pageanchor=false}%
\begin{titlepage}
\thispagestyle{empty}
\vspace*{0.25cm}
{\color{accent}\large ORDINAL ARITHMETIC AND SURREAL NUMBERS\par}
\vspace{0.4cm}
{\color{ink}\huge\bfseries Gonshor's Product-Birthday Conjecture\par}
\vspace{0.45cm}
{\color{ink}\LARGE An exact formula for ordinal-support series\par}
\vspace{0.45cm}
{\large A partial solution, an equality classification,\par
and exact reproducible checks\par}
\vspace{0.45cm}
{Research draft prepared with ChatGPT\par
20 September 2026\par}
\vspace{0.45cm}
\begin{resultbox}
\textbf{Scope of the result.}
The conjectured inequality
\[
 \bd(xy)\leq \bd(x)\np\bd(y)
\]
is proved here for all surreal normal forms with real coefficients and
exponents that are negatives of ordinals, allowing arbitrary set-sized
transfinite support. The proof also identifies every equality case in this
ring. The conjecture for arbitrary surreal numbers is \textbf{not settled}
by this article.
\end{resultbox}
\vspace{0.15cm}
\begin{abstract}
Gonshor's product-birthday conjecture asks whether the canonical sign length
of a surreal product is bounded by the natural product of the sign lengths
of its factors. We investigate the ring
\(\A=\{\sum_{\alpha\in S}c_\alpha\omega^{-\alpha}\}\), where
\(S\) is any set of ordinals and the nonzero coefficients are real.
Specializing the classical reduced-sign-expansion theorem gives an exact
birthday formula determined by the supremum of the support, the last
coefficient when one exists, and one endpoint correction. This yields a
sharp support bound for products and shows that multiplication does not
increase the leading Cantor degree of the birthday inside \(\A\).
Consequently the conjecture holds throughout \(\A\), strictly except for
zero factors, factors equal to \(\pm1\), and pairs of ordinary integers.
We also obtain a four-coefficient formula for finite-support product
birthdays and describe natural set-sized local subrings. Exact rational
computations check the endpoint formula on 64,401 series and the product
results on 96,636 pairs. The transfinite statements are established by
written proofs, not by these finite tests. Historical priority for these
particular special-case results has not been established.
\end{abstract}
\vfill
{\small\color{ink}Companion materials: Python source, deterministic verification
results, bibliography, and build instructions. No proof-assistant
verification is claimed.}
\end{titlepage}
\hypersetup{pageanchor=true}%

\tableofcontents
\newpage

\section{The problem, its status, and the outcome}

\subsection{The conjecture actually being investigated}
A surreal number has a canonical sign expansion of ordinal length. Write
\(\bd(x)\) for that length, with \(\bd(0)=0\). The question is
\begin{equation}\label{eq:conjecture}
 \boxed{\quad \bd(xy)\leq\bd(x)\np\bd(y)
          \qquad(x,y\in\No).\quad}
\end{equation}
Here \(\np\) denotes the \emph{natural}, or Hessenberg, product of ordinals.
The corresponding additive bound uses natural sum:
\(\bd(x+y)\leq\bd(x)\ns\bd(y)\).

The Open Problem Garden entry attributes \eqref{eq:conjecture} to
Gonshor's 1986 book \cite{Gonshor,OPG}. More substantially, in an
August 2024 account Philip Ehrlich identified the question as Gonshor's
conjecture and reported the monomial case proved with van den Dries:
\(x=r\omega^a\), \(y=s\omega^b\), with arbitrary surreal \(a,b\)
and real \(r,s\). His account did not report a general solution
\cite{EhrlichMO}. No general resolution was located in the sources
consulted for this draft on 20 September 2026. That is a statement about
the literature search, not a certification that every relevant publication
has been examined.

There is an elementary reason that the ordinal operations must be made
explicit. With \(x=\omega+1\) and \(y=2\), their surreal product is the
ordinal \(\omega\cdot2+2\). Its birthday is therefore
\(\omega\cdot2+2\), whereas the \emph{ordinary} ordinal product
\((\omega+1)\cdot2=\omega\cdot2+1\) is smaller. The natural product
has the correct value \((\omega+1)\np2=\omega\cdot2+2\).

\subsection{Why recursive expression rank is not the answer}
The rank of an unreduced recursive expression for a surreal number can
exceed the canonical sign length of the number it represents. For example,
Roughan's calculation of the recursive product form for \(2\times3\)
gives generation 12 \cite{Roughan}. The value of that form is the integer
6, whose canonical birthday is 6. Thus the form calculation does not
contradict \eqref{eq:conjecture}. In this article every occurrence of
\(\bd\) refers to the canonical number, never to an expression tree.

This distinction matters for both research and implementation. Expanding
Conway's recursive multiplication rule and measuring its depth does not
by itself measure the birthday of its output value. Our computations
instead evaluate a canonical normal-form sign-length rule.

\subsection{The result and its logical dependencies}
The main domain is
\begin{equation}\label{eq:ring-intro}
 \A=\{0\}\cup
 \left\{\sum_{\alpha\in S}c_\alpha\omega^{-\alpha}:
 S\subseteq\On\text{ is a nonempty set},\quad
 c_\alpha\in\R\setminus\{0\}\right\}.
\end{equation}
The sum is a Conway normal form, ordered by increasing \(\alpha\).
These exponents are \emph{negative ordinals as surreal numbers}; they are
not arbitrary negative surreal exponents. This ring contains \(\R\),
\(\R[[\omega^{-1}]]\), and many series with uncountable or otherwise
transfinite support. It does not contain \(\omega\), or even
\(\omega^{-1/2}\).

The route to the main theorem is
\[
\begin{gathered}
 \text{classical reduced sign expansions}
 \ \Longrightarrow\ \text{exact endpoint formula}\\
 \Longrightarrow\ \text{support and degree bounds}
 \ \Longrightarrow\ \eqref{eq:conjecture}\text{ on }\A.
\end{gathered}
\]
The endpoint formula and the deductions from it are proved below.
The substantial imported ingredient is the classical sign-expansion
theorem. Accordingly, this is not a foundational reconstruction of all
surreal arithmetic. It is a self-contained derivation of a special-case
product theorem from an explicitly stated standard ingredient.

\begin{resultbox}
\textbf{Research status.}
The general conjecture remains outside the conclusion of this draft.
The results below are proposed as rigorous special-case theorems, with
complete arguments relative to the classical normal-form sign theorem.
They have not been peer reviewed or formally checked, and no claim of
historical priority is made.
\end{resultbox}

\section{Conventions and elementary ordinal facts}

\subsection{Two different kinds of ordinal arithmetic}
We use \(\N=\{0,1,2,\ldots\}\). A finite birthday means a finite
ordinal sign length, not merely a surreal value bounded by an integer.
Every element of \(\A\) has finite magnitude; finite birthday
characterizes its dyadic real constants.

Ordinary ordinal addition and multiplication are written \(\alpha+\beta\)
and \(\alpha\cdot\beta\). They describe ordered concatenation and its
iteration. The operations \(\ns\) and \(\np\) are natural sum and
natural product. Surreal addition and multiplication use the usual symbols;
when both operands are ordinals embedded in \(\No\), they agree with the
natural operations \cite[Section~6]{Berarducci}.

In particular, throughout this paper \(\omega\cdot\alpha\) in a
\emph{birthday formula} is an ordinary ordinal product. One must not replace
it by \(\omega\np\alpha\). Likewise, \(1+\rho\) in an exponent below
is ordinary addition and can differ from \(\rho+1\).

If a common decreasing list of Cantor exponents is chosen, then
\[
 \alpha=\sum_{j=1}^m\omega^{\rho_j}a_j,
 \qquad
 \beta=\sum_{j=1}^m\omega^{\rho_j}b_j,
 \qquad \rho_1>\cdots>\rho_m,
\]
where zero coefficients are allowed for alignment, and
\begin{equation}\label{eq:natural-sum}
 \alpha\ns\beta=\sum_{j=1}^m\omega^{\rho_j}(a_j+b_j).
\end{equation}
Natural product is polynomial multiplication, with exponent addition
itself interpreted naturally:
\begin{equation}\label{eq:natural-product}
 \left(\sum_i\omega^{\xi_i}a_i\right)
 \np\left(\sum_j\omega^{\eta_j}b_j\right)
 =\mathop{\bigoplus}_{i,j}\omega^{\xi_i\ns\eta_j}a_i b_j.
\end{equation}
These formulas also fix the notation for arbitrarily large ordinals.

\begin{definition}
For a nonzero ordinal \(\theta\), write \(\deg\theta\) for its leading
Cantor exponent and \(\lc\theta\) for its leading positive integer
coefficient. Thus
\[
 \theta=\omega^{\deg\theta}\cdot\lc\theta+\theta',
 \qquad \theta'<\omega^{\deg\theta}.
\]
We leave \(\deg0\) undefined. In particular, a positive finite ordinal
has degree zero.
\end{definition}

\begin{lemma}\label[lemma]{lem:degree-arithmetic}
For nonzero ordinals \(\theta,\eta\),
\begin{align}
 \deg(\theta\ns\eta)&=\max(\deg\theta,\deg\eta),\label{eq:deg-sum}\\
 \deg(\theta\np\eta)&=\deg\theta\ns\deg\eta,\label{eq:deg-product}\\
 \deg(\omega\cdot\theta)&=1+\deg\theta.\label{eq:deg-left-omega}
\end{align}
If \(d,e>0\), then
\begin{equation}\label{eq:degree-gap}
 d\ns e\geq\max(d,e)+1.
\end{equation}
Natural sum is strictly increasing in each argument and is therefore
monotone in both arguments jointly.
\end{lemma}
\begin{proof}
Equation \eqref{eq:deg-sum} follows by aligning Cantor normal forms.
In \eqref{eq:natural-product}, the largest exponent is the natural sum
of the largest exponent from each factor, proving
\eqref{eq:deg-product}. Ordinary left multiplication by \(\omega\)
changes a leading term \(\omega^d c\) to \(\omega^{1+d}c\), proving
\eqref{eq:deg-left-omega}.

For strict monotonicity, compare two ordinal Cantor normal forms at
the greatest exponent where their coefficients differ. Adding the same
coefficient list does not change that first comparison. Thus
\(a<a'\) implies \(a\ns b<a'\ns b\). Taking \(b>0\) shows that
\(a\ns b>a\); interchanging the arguments gives
\(a\ns b>b\). For ordinals, being strictly greater than \(D\) means
being at least \(D+1\), proving \eqref{eq:degree-gap}.
\end{proof}

\begin{lemma}\label[lemma]{lem:interval}
If \(\delta\leq\alpha\), the interval \([\delta,\alpha)\), in its
inherited order, has an order type \(\eta\) satisfying
\(\delta+\eta=\alpha\). For \(\eta>0\), deleting its first point
leaves an order type \(\eta'\) satisfying \(1+\eta'=\eta\).
Furthermore,
\[
 \omega\cdot\delta+\omega\cdot\eta=\omega\cdot\alpha,
 \qquad
 \sup_i\omega\cdot\alpha_i
   =\omega\cdot\sup_i\alpha_i
\]
for an increasing nonempty family with no last index.
\end{lemma}
\begin{proof}
The ordinal \(\alpha\) is the concatenation of its initial segment
\(\delta\) and the stated interval. The deletion statement is the same
observation with initial segment 1. The identities follow respectively
from distributivity of left multiplication over ordinary ordinal addition
and from the limit clause in the definition of ordinal multiplication.
\end{proof}

\section{The ordinal-support ring}

\subsection{Normal forms and the meaning of the sums}
We use Conway's normal-form representation: every surreal has a unique
set-sized series \(\sum_{i<\nu}r_i\omega^{a_i}\) with nonzero real
coefficients and strictly decreasing surreal exponents. Addition and
multiplication are the corresponding series operations, and
\(\omega^a\omega^b=\omega^{a+b}\)
\cite[Sections~8--10]{Berarducci}. The notation \(\omega^a\) here
means Conway's omega-map, not a substitution into a separately defined
real exponential function.

For \(x\in\A\), let \(S(x)\) be its \emph{ordinal-index support}:
\[
 S(x)=\{\alpha:c_\alpha\ne0\}.
\]
This is a set of nonnegative ordinals, not the usual support of actual
exponents \(\{-\alpha:\alpha\in S(x)\}\). A subset of the ordinals
is automatically well ordered. Its increasing enumeration makes
\(-\alpha\) a decreasing sequence, so every sum in
\eqref{eq:ring-intro} is a legitimate normal form. No analytic convergence
condition or boundedness condition on the real coefficients is imposed.
In the common Hahn notation one can write \(t^\alpha=\omega^{-\alpha}\).
Here the index addition in \(t^\alpha t^\beta=t^{\alpha\ns\beta}\)
is natural sum: the ordinal indices form a monoid, rather than the whole
additive exponent group of a Hahn field.

\begin{proposition}\label[proposition]{prop:closure}
The class \(\A\) is a commutative subring of \(\No\). If
\(x=\sum c_\alpha\omega^{-\alpha}\) and
\(y=\sum d_\beta\omega^{-\beta}\), then
\begin{equation}\label{eq:convolution}
 xy=\sum_\gamma p_\gamma\omega^{-\gamma},
 \qquad
 p_\gamma=\sum_{\alpha\ns\beta=\gamma}c_\alpha d_\beta,
\end{equation}
with absent coefficients interpreted as zero. Each inner sum is finite.
\end{proposition}
\begin{proof}
Surreal addition on embedded ordinals is natural sum. Therefore
\[
 (-\alpha)+(-\beta)=-(\alpha\ns\beta),\qquad
 \omega^{-\alpha}\omega^{-\beta}
       =\omega^{-(\alpha\ns\beta)}.
\]
To check local finiteness directly, write
\(\gamma=\sum_{j=1}^m\omega^{\rho_j}n_j\) in Cantor normal form.
A pair of ordinals with natural sum \(\gamma\) is obtained by splitting
each integer coefficient \(n_j\) between its two components. No new
Cantor exponent can appear, since the coefficients are nonnegative.
Consequently there are exactly
\begin{equation}\label{eq:fiber-count}
 \#\{(\alpha,\beta):\alpha\ns\beta=\gamma\}
     =\prod_{j=1}^m(n_j+1)
\end{equation}
ordered pairs. For \(\gamma=0\), the empty product is 1.

The candidate product indices form the set
\(\{\alpha\ns\beta:\alpha\in S(x),\beta\in S(y)\}\).
It is well ordered, and every coefficient is a finite sum. Thus the normal
form product exists and belongs to \(\A\). Addition and negation are
coefficientwise and preserve membership. The real constants give the
identity element. All ring laws are inherited from \(\No\).
\end{proof}

In particular, \(\A\) has no zero divisors. It is not a subfield:
\(\omega^{-1}\in\A\), whereas its inverse \(\omega\) is not in
\(\A\). This failure is a genuine restriction of the result, not just a
matter of terminology.

\section{The classical sign rule in the all-minus case}

\subsection{Real coefficients}
Put \(\D=\Z[1/2]\). A surreal has finite birthday exactly when it is a
dyadic real. A non-dyadic real has birthday \(\omega\)
\cite[Section~5]{Berarducci}. For \(0\ne c=m/2^k\in\D\) in reduced
form, \(k\geq0\), the finite birthday is
\begin{equation}\label{eq:dyadic-birthday}
 d(c):=\bd(c)=\lceil |c|\rceil+k.
\end{equation}
For integer \(c\), this says \(d(c)=|c|\). The formula is also recorded
as Lemma~1 of \cite{Roughan}; it follows by the successive bisections in
the dyadic sign tree.

For use in a sign block, define the coefficient-tail length
\begin{equation}\label{eq:tail}
 \tau(c)=
 \begin{cases}
 d(c)-1,&c\in\D\setminus\{0\},\\
 \omega,&c\in\R\setminus\D.
 \end{cases}
\end{equation}
Deleting the initial sign of a finite sign string shortens it by one.
Deleting the initial sign of an \(\omega\)-long real sign string leaves
length \(\omega\). Reversing all signs, as needed for a negative
coefficient, does not change any of these lengths.

\subsection{Precisely the imported ingredient}
We use Gonshor's reduced-sign-expansion theorem, in the formulation of
Bournez and Guilmant, Definition~2.20 and Theorem~2.21
\cite{BournezGuilmant}; their theorem cites Gonshor's
Theorems~5.11--5.12. Its needed specialization is as follows.

\begin{resultbox}
For successive normal-form exponents \(-\alpha_i\), begin with their
all-minus sign strings. Delete each minus position already shared with
an earlier exponent. At a successor term, delete one additional minus
after the preceding exponent if its coefficient is non-dyadic. Start each term with one initial plus, then append an \(\omega\)-long
block for each remaining minus. Append the sign string of the absolute
coefficient without its initial plus. Reverse the whole term for a negative
coefficient. Concatenate the term blocks in normal-form order.
\end{resultbox}

Only the lengths of these blocks are needed. The specialization is
particularly simple because \(-\alpha\) has \(\alpha\) minus signs and
no plus signs. In the general sign theorem, plus signs in the exponent
change the sizes of subsequent blocks; that complication is absent here.

\subsection{A usable block-length formula}
Write the support in increasing order as
\(S=\{\alpha_i:i<\nu\}\), with coefficient \(c_i\), and put
\[
 \delta_i=\sup_{j<i}\alpha_j,\qquad\sup\varnothing=0.
\]
The positions shared with earlier all-minus exponents are exactly
\([0,\delta_i)\): a position \(\xi\) is shared precisely when
\(\xi<\alpha_j\) for some \(j<i\). Therefore the remaining interval
is \([\delta_i,\alpha_i)\). Let its order type be \(\eta_i\).

When \(i\) is a successor and \(c_{i-1}\notin\D\), delete the first
point of this interval and call the resulting order type \(\eta_i'\).
The interval is nonempty in that case: \(\delta_i=\alpha_{i-1}<\alpha_i\),
and the extra deleted position is exactly \(\alpha_{i-1}\). There is
no preceding coefficient at a nonzero limit position, so only shared
minus signs are deleted there. In all other cases put
\(\eta_i'=\eta_i\). The specialized sign rule gives
\begin{equation}\label{eq:block-formula}
 B_i=1+\omega\cdot\eta_i'+\tau(c_i),
 \qquad
 \bd(x)=\sum_{i<\nu} B_i,
\end{equation}
where every operation on the right is ordinary ordinal arithmetic.

This equation is already a transfinite algorithm, but it appears to
require the whole support. The next theorem compresses it to endpoint
information. The point is not that earlier signs literally disappear.
Rather, finite initial lengths are absorbed when later \(\omega\)-blocks
are appended, as in \(7+\omega=\omega\).

\section{An exact endpoint formula}\label{sec:endpoint}

\begin{theorem}[Exact birthday]\label[theorem]{thm:endpoint}
Let \(0\ne x=\sum_{\alpha\in S}c_\alpha\omega^{-\alpha}\in\A\)
and put \(\sigma=\sup S\).

If \(S\) has no greatest element, then
\begin{equation}\label{eq:no-maximum}
 \boxed{\bd(x)=\omega\cdot\sigma.}
\end{equation}
If \(S\) has greatest element \(a\) and \(c_a\notin\D\), then
\begin{equation}\label{eq:nondyadic-last}
 \boxed{\bd(x)=\omega\cdot(a+1).}
\end{equation}
If \(S\) has greatest element \(a\) and \(c_a\in\D\), put
\(n=d(c_a)\) and define \(e(x)\in\{0,1\}\) by
\begin{equation}\label{eq:correction}
 e(x)=1\quad\Longleftrightarrow\quad
 \begin{cases}
 \sup(S\setminus\{a\})=a,\quad\text{or}\\
 a=b+1\text{ for some }b\in S\text{ with }c_b\notin\D.
 \end{cases}
\end{equation}
Here \(\sup\varnothing=0\), so the first condition includes the
constant case \(a=0\). Then
\begin{equation}\label{eq:dyadic-last}
 \boxed{\bd(x)=\omega\cdot a+(n-1+e(x)).}
\end{equation}
\end{theorem}

There are two different reasons for the correction. The last exponent
may be a limit of preceding support indices, leaving no new exponent
signs; or it may be the immediate ordinal successor of a preceding index
with a non-dyadic coefficient, causing the one new minus to be deleted.
These mechanisms must not be conflated.

\begin{proof}
Enumerate \(S\) as in \eqref{eq:block-formula}. We prove the statements
for all initial normal forms by transfinite induction. For a segment with
last index \(\alpha_i\), let \(P_i\) denote its sign length. Along the
way the following bounds hold:
\begin{equation}\label{eq:prefix-bounds}
 \omega\cdot\alpha_i\leq P_i
       \leq\omega\cdot(\alpha_i+1).
\end{equation}
They are immediate from the asserted formulas once each stage is proved.

\paragraph{The first term.}
If \(\alpha_0=0\), the exponent string is empty. Equation
\eqref{eq:block-formula} gives \(1+\tau(c_0)\), which is \(d(c_0)\)
for a dyadic coefficient and \(\omega\) otherwise. This is precisely
\eqref{eq:dyadic-last} with \(e=1\), or
\eqref{eq:nondyadic-last}.

If \(\alpha_0>0\), no minus is deleted and the block has length
\(1+\omega\cdot\alpha_0+\tau(c_0)\). The initial 1 is absorbed by the
nonzero \(\omega\)-multiple. We obtain
\(\omega\cdot\alpha_0+(d(c_0)-1)\) in the dyadic case, with \(e=0\),
and \(\omega\cdot(\alpha_0+1)\) in the non-dyadic case.

\paragraph{A successor term after a dyadic coefficient.}
Suppose the preceding index is \(\beta\), the new one is \(\alpha>\beta\),
and the preceding coefficient is dyadic. Its initial normal form has
length \(\omega\cdot\beta+t\) for some finite \(t\), by induction.
Let \(\eta=\otp[\beta,\alpha)\), so that
\(\beta+\eta=\alpha\) and \(\eta>0\). No extra minus is deleted.
Appending the new block gives
\begin{align*}
 (\omega\cdot\beta+t)+(1+\omega\cdot\eta+\tau(c_\alpha))
 &=\omega\cdot\beta+\omega\cdot\eta+\tau(c_\alpha)\\
 &=\omega\cdot\alpha+\tau(c_\alpha).
\end{align*}
The first equality uses that a finite ordinal followed by a nonzero
\(\omega\)-multiple is absorbed. This proves the formulas, with no
endpoint correction in the dyadic last-coefficient case.

\paragraph{A successor term after a non-dyadic coefficient.}
The preceding segment has length \(\omega\cdot(\beta+1)\).
Delete the first point of \([\beta,\alpha)\), leaving order type
\(\eta'\). Thus \(1+\eta'=\eta\) and
\(\beta+1+\eta'=\alpha\).
If \(\eta'>0\), appending the new block gives
\[
 \omega\cdot(\beta+1)+(1+\omega\cdot\eta'+\tau(c_\alpha))
   =\omega\cdot\alpha+\tau(c_\alpha).
\]
Again the asserted formulas follow without correction.

If \(\eta'=0\), then \(\alpha=\beta+1\). There are no retained
minus signs. Appending the new block gives
\[
 \omega\cdot\alpha+(1+\tau(c_\alpha)).
\]
For a dyadic last coefficient this is
\(\omega\cdot\alpha+d(c_\alpha)\), giving exactly \(e=1\).
For a non-dyadic last coefficient it is
\(\omega\cdot(\alpha+1)\), as required.

\paragraph{A term at a nonzero limit position.}
Suppose \(i\) is a nonzero limit ordinal and let
\(\delta=\sup_{j<i}\alpha_j\). The concatenation preceding term \(i\)
has length \(\sup_{j<i}P_j\). The bounds
\eqref{eq:prefix-bounds} and ordinary ordinal continuity give
\begin{equation}\label{eq:limit-prefix}
 \sup_{j<i}P_j=\omega\cdot\delta.
\end{equation}
Indeed the preceding indices have no maximum and are cofinal in
\(\delta\). For each \(j<i\), also \(j+1<i\), and strict increase gives
\(\alpha_j+1\leq\alpha_{j+1}\leq\delta\). Hence
\(\sup_{j<i}(\alpha_j+1)=\sup_{j<i}\alpha_j=\delta\), and both
bounding suprema are \(\omega\cdot\delta\). This argument uses the
whole cofinal support and does not require it to be countable.

There is no additional successor deletion at a limit position. If
\(\delta<\alpha_i\), the retained interval is nonempty. The same
absorption and concatenation calculation as above gives
\(\omega\cdot\alpha_i+\tau(c_i)\). If \(\delta=\alpha_i\), the
retained interval is empty, so the result is
\(\omega\cdot\alpha_i+(1+\tau(c_i))\). The latter is the cofinal
endpoint correction in \eqref{eq:correction}.

\paragraph{A complete normal form with no last term.}
Its support has no maximum. Apply the same squeeze argument to all
initial segments ending in a term. Equation \eqref{eq:limit-prefix}, now
with \(\delta=\sigma\), gives \eqref{eq:no-maximum}.
All steps preserve \eqref{eq:prefix-bounds}, completing the induction.
\end{proof}

\begin{corollary}\label[corollary]{cor:support-sandwich}
For \(0\ne x\in\A\), with \(\sigma=\sup S(x)\),
\begin{equation}\label{eq:support-sandwich}
 \omega\cdot\sigma\leq\bd(x)
       \leq\omega\cdot(\sigma+1).
\end{equation}
If \(\bd(x)\) is infinite, then
\begin{equation}\label{eq:support-degree}
 \deg\bd(x)=1+\deg\bigl(\max\{1,\sigma\}\bigr).
\end{equation}
\end{corollary}
\begin{proof}
The bounds follow term by term from \cref{thm:endpoint}. If \(\sigma>0\),
the bounding ordinals have the same leading Cantor exponent, since
\(\deg(\sigma+1)=\deg\sigma\). This gives
\(\deg\bd(x)=1+\deg\sigma\). If \(\sigma=0\), the number is a real
constant; an infinite birthday then means a non-dyadic constant and
\(\bd(x)=\omega\), which also gives \eqref{eq:support-degree}.
\end{proof}

\section{Products: a support estimate and a degree gap}

\begin{theorem}[Sharp support bound]\label[theorem]{thm:support-product}
Let \(0\ne x,y\in\A\), and let
\(\alpha=\sup S(x)\), \(\beta=\sup S(y)\). Then
\begin{equation}\label{eq:support-product}
 \boxed{\bd(xy)\leq\omega\cdot\bigl((\alpha\ns\beta)+1\bigr).}
\end{equation}
For every prescribed pair \(\alpha,\beta\), equality is attainable by
some factors with those support suprema.
\end{theorem}
\begin{proof}
Every surviving product index is \(a\ns b\) with \(a\leq\alpha\)
and \(b\leq\beta\). Monotonicity gives
\(a\ns b\leq\alpha\ns\beta\), so
\(\sup S(xy)\leq\alpha\ns\beta\). Now apply
\eqref{eq:support-sandwich}, which applies because the nonzero factors
have nonzero product in \(\No\). Coefficient cancellation only removes
indices and cannot spoil this upper bound.

For sharpness choose
\(x=\frac13\omega^{-\alpha}\) and
\(y=\omega^{-\beta}\). Their product is the monomial
\(\frac13\omega^{-(\alpha\ns\beta)}\), whose birthday is the
right-hand side of \eqref{eq:support-product} by
\eqref{eq:nondyadic-last}.
\end{proof}

The sharpness assertion fixes the support suprema, not the two input
birthdays. These are different optimization problems. In particular it
does not imply that the conjectured birthday-product bound is sharp for
general non-real factors in this ring.

\begin{proposition}[Strict degree gap]\label[proposition]{prop:infinite-gap}
Suppose \(x,y\in\A\) have infinite birthdays. Put
\(d=\deg\bd(x)\), \(e=\deg\bd(y)\), and \(D=\max(d,e)\). Then
\begin{equation}\label{eq:gap}
 \bd(xy)<\omega^{D+1}
       \leq\omega^{d\ns e}
       \leq\bd(x)\np\bd(y).
\end{equation}
In particular the conjectured inequality is strict for this pair.
\end{proposition}
\begin{proof}
The factors and their product are nonzero. Write
\(\alpha=\sup S(x)\), \(\beta=\sup S(y)\) and set
\[
 r=\max\bigl(\deg\max\{1,\alpha\},
                  \deg\max\{1,\beta\}\bigr).
\]
The ordinal \((\alpha\ns\beta)+1\) has leading exponent \(r\).
This is immediate by aligning Cantor forms, including the case
\(\alpha=\beta=0\). Thus the upper bound in
\eqref{eq:support-product} has degree \(1+r\).
On the other hand, \eqref{eq:support-degree} gives
\[
 D=\max\bigl(1+\deg\max\{1,\alpha\},
              1+\deg\max\{1,\beta\}\bigr)=1+r.
\]
It follows that \(\deg\bd(xy)\leq D\). Every ordinal of degree at
most \(D\) is strictly below \(\omega^{D+1}\), proving the first
inequality in \eqref{eq:gap}.

Both \(d,e\) are positive because both input birthdays are infinite.
Hence \cref{lem:degree-arithmetic} gives \(d\ns e\geq D+1\).
Finally \(\bd(x)\np\bd(y)\) has leading exponent \(d\ns e\) and
positive leading coefficient, so it is at least \(\omega^{d\ns e}\).
\end{proof}

\section{The product conjecture and all equality cases on the ring}

\subsection{Finite factors}
\begin{lemma}[Dyadic factors]\label[lemma]{lem:dyadic-product}
For dyadic reals \(u,v\),
\(\bd(uv)\leq\bd(u)\bd(v)\). Equality holds precisely when a factor
is zero, a factor is \(\pm1\), or both factors are integers.
\end{lemma}
\begin{proof}
The zero cases are immediate. For nonzero \(u,v\), put
\(a=\lceil|u|\rceil\), \(b=\lceil|v|\rceil\), and let \(k,l\)
be the respective reduced denominator exponents. Thus
\(\bd(u)=a+k\) and \(\bd(v)=b+l\). The product has denominator
exponent at most \(k+l\), and \(\lceil|uv|\rceil\leq ab\). Hence
\begin{equation}\label{eq:dyadic-comparison}
 \bd(uv)\leq ab+k+l\leq(a+k)(b+l).
\end{equation}
The difference in the second inequality is
\[
 (a+k)(b+l)-(ab+k+l)=l(a-1)+k(b-1)+kl.
\]
If \(k,l>0\), this is positive. If \(k>0,l=0\), it is positive
unless \(b=1\), in which case the integer \(v\) is \(\pm1\).
The remaining mixed case is symmetric. If \(k=l=0\), both numbers
are integers and equality is immediate. Multiplication by \(\pm1\)
also preserves birthday, completing the classification.
\end{proof}

\begin{lemma}[A dyadic scalar cannot increase the leading birthday data]
\label[lemma]{lem:scalar}
Let \(0\ne q\in\D\) and \(x\in\A\) have infinite birthday. The
ordered pair
\(\bigl(\deg\bd(qx),\lc\bd(qx)\bigr)\) is lexicographically no larger
than \(\bigl(\deg\bd(x),\lc\bd(x)\bigr)\). Consequently, if
\(q\notin\{1,-1\}\), then
\begin{equation}\label{eq:scalar-strict}
 \bd(qx)<\bd(q)\np\bd(x).
\end{equation}
\end{lemma}
\begin{proof}
Multiplication by the nonzero constant \(q\) preserves the index support.
If that support has no maximum, \eqref{eq:no-maximum} gives exactly the
same birthday.

Suppose it has maximum \(a\). A dyadic coefficient remains dyadic after
multiplication by \(q\). If the last coefficient is dyadic, both birthdays
have the form \(\omega\cdot a+\text{finite}\). Since \(\bd(x)\) is
infinite, \(a>0\), and the finite tails do not change the leading Cantor
data. The endpoint correction may change, but only within those finite
tails.

If the last coefficient is non-dyadic, it may remain non-dyadic or become
dyadic. In the former case both birthdays equal
\(\omega\cdot(a+1)\). In the latter the new birthday is
\(\omega\cdot a+\text{finite}<\omega\cdot(a+1)\), so its leading
Cantor data cannot increase. This includes \(a=0\), when the new
birthday may be finite.

Finally \(m=\bd(q)\) is a positive integer, and \(m\geq2\) unless
\(q=\pm1\). Natural multiplication by \(m\) preserves the leading
exponent of \(\bd(x)\) and multiplies its leading coefficient by \(m\).
The preceding leading-data comparison makes \eqref{eq:scalar-strict}
strict: a smaller leading exponent already gives a smaller ordinal; at
equal leading exponent, the new leading coefficient is at most the old
one and therefore strictly below its multiple by \(m\geq2\).
\end{proof}

\begin{remark}
One must not claim that a non-dyadic real coefficient remains non-dyadic
under a dyadic scalar: \(3\cdot\frac13=1\). Nor is the entire birthday
necessarily nonincreasing. For example,
\[
 \bd\left(\frac13+\omega^{-1}\right)=\omega+1,
 \qquad
 \bd\left(3\left(\frac13+\omega^{-1}\right)\right)=\omega+2.
\]
What is controlled in \cref{lem:scalar} is the leading Cantor exponent and
coefficient, not the finite tail.
\end{remark}

\subsection{Main theorem}
\begin{theorem}[Product birthdays on \(\A\)]\label[theorem]{thm:main}
For every \(x,y\in\A\),
\begin{equation}\label{eq:main}
 \boxed{\bd(xy)\leq\bd(x)\np\bd(y).}
\end{equation}
Equality holds if and only if at least one of the following holds:
\begin{enumerate}[label=(\roman*),topsep=2pt,itemsep=1pt]
 \item \(x=0\) or \(y=0\);
 \item \(x\in\{1,-1\}\) or \(y\in\{1,-1\}\);
 \item \(x\) and \(y\) are both ordinary integers.
\end{enumerate}
\end{theorem}
\begin{proof}
Zero factors are immediate. If both nonzero factors have finite birthdays,
they are dyadic reals and \cref{lem:dyadic-product} gives the inequality
and equality classification. If exactly one factor has finite birthday,
it is a nonzero dyadic constant. Multiplication by \(\pm1\) gives
equality; every other such constant gives strict inequality by
\cref{lem:scalar}. If both birthdays are infinite,
\cref{prop:infinite-gap} gives strict inequality. These alternatives
exhaust all pairs in \(\A\), and each displayed equality case does
indeed give equality.
\end{proof}

\begin{corollary}[Degree does not rise]\label[corollary]{cor:degree}
For nonzero \(x,y\in\A\),
\begin{equation}\label{eq:degree-nonincreasing}
 \deg\bd(xy)\leq\max\bigl(\deg\bd(x),\deg\bd(y)\bigr).
\end{equation}
Thus for every positive integer \(n\),
\(\deg\bd(x^n)\leq\deg\bd(x)\).
\end{corollary}
\begin{proof}
When both birthdays are infinite this was proved in
\cref{prop:infinite-gap}. With one finite birthday use
\cref{lem:scalar}; with both finite, all the degrees are zero.
The power statement follows by induction.
\end{proof}

The degree conclusion is much stronger than the degree estimate that
would follow merely from \eqref{eq:main}. That weaker estimate permits
\(\deg\bd(x)\ns\deg\bd(y)\), whereas here only their maximum is
needed. This explains why equality is so rare within the ring.

\section{A four-coefficient formula for finite supports}\label{sec:four}

Suppose \(x,y\in\A\setminus\{0\}\) have finite index supports, with
maxima \(\alpha,\beta\). These maxima pick out the last normal-form
terms. Put
\[
 \gamma=\alpha\ns\beta,\qquad p=c_\alpha d_\beta.
\]
For an ordinal \(\theta\), let \(\theta^-\) denote its immediate
predecessor when it exists; it is undefined at zero and at a limit.
An absent coefficient, including a coefficient at an undefined
predecessor, is interpreted as zero. Define
\begin{equation}\label{eq:near-top}
 q=c_{\alpha^-}d_\beta+c_\alpha d_{\beta^-}.
\end{equation}

\begin{proposition}[No full convolution is needed]\label[proposition]{prop:four}
With the notation above, the product birthday is
\begin{equation}\label{eq:four}
 \bd(xy)=
 \begin{cases}
 \omega\cdot(\gamma+1),&p\notin\D,\\
 \omega\cdot\gamma+(d(p)-1+E),&p\in\D,
 \end{cases}
\end{equation}
where
\[
 E=1\quad\Longleftrightarrow\quad
 \gamma=0\quad\text{or}\quad
 (\gamma\text{ is a successor ordinal and }q\notin\D).
\]
After locating the two support maxima, at most the top and predecessor
coefficient of each factor is needed: four coefficient lookups in total. The full product support need not be constructed.
\end{proposition}
\begin{proof}
Strict monotonicity of natural sum shows that any pair of indices other
than \((\alpha,\beta)\) has sum strictly less than \(\gamma\).
Thus the product has maximum index exactly \(\gamma\), with nonzero
coefficient \(p\): there is only one contribution there, so no
cancellation can remove that endpoint.

If \(\gamma\) is a successor, a decomposition
\(a\ns b=\gamma^-\), with \(a\leq\alpha\), \(b\leq\beta\), can
only be \((\alpha^-,\beta)\) or \((\alpha,\beta^-)\), when defined.
To see this, align the Cantor forms of \(a,b,\alpha,\beta\), inserting
zero coefficients as necessary. The target \(\gamma^-\) differs from
\(\gamma\) only by subtracting one from the constant coefficient.
Suppose some positive-exponent coefficient of an input differs from that
of its corresponding maximum, and choose the greatest exponent \(\xi>0\)
where either input differs. Above \(\xi\) each input agrees with its
maximum. The ordinal bounds \(a\leq\alpha\) and \(b\leq\beta\)
therefore force both coefficients at \(\xi\) to be no larger, with at
least one strictly smaller. Their sum cannot equal the coefficient of
\(\gamma^-\) at \(\xi\), a contradiction. This first-difference
argument is needed because ordinal order is lexicographic on Cantor
coefficients.

Thus all positive-exponent coefficients are unchanged. Each constant
coefficient can only decrease, and their total decreases by exactly one.
Exactly one input is its maximum's immediate predecessor and the other
is unchanged. The coefficient at
\(\gamma^-\) is consequently \(q\), including the possibility that
it cancels to zero; zero is dyadic and gives no successor correction.

The product support is finite, so it cannot accumulate at a positive
last index. Its endpoint correction is therefore precisely the one in
the statement. Apply \cref{thm:endpoint}.
\end{proof}

\begin{remark}
The word ``four'' counts coefficient lookups, not bit operations.
Ordinal-index arithmetic and exact coefficient multiplication still have
costs, and locating a maximum in an unsorted input requires a scan.
The claim is that a quadratic-size convolution can be avoided when only
the birthday is requested. For arbitrary real coefficients, the formula
requires exact decisions about \(p\in\D\) and \(q\in\D\); it does not
supply such decisions from computable-real names. The executable rational
case decides dyadic status by inspecting reduced denominators. For infinite
supports with a last limit index,
the cofinal endpoint correction must also be considered; the proposition
is deliberately restricted to finite supports.
\end{remark}

For example, with \(t=\omega^{-1}\), take
\(x=\tfrac13+t\) and \(y=-\tfrac13+t\). Their top indices are both
\(1\), their top coefficient product is \(p=1\), and their predecessor
coefficient is \(q=\tfrac13-\tfrac13=0\). Thus
\[
 xy=-\tfrac19+t^2,\qquad \bd(xy)=\omega\cdot2.
\]
The non-dyadic constant of the product is not adjacent to its top index,
so it creates no correction there. Replacing \(y\) by \(x\) instead
gives \(q=2/3\) and \(\bd(x^2)=\omega\cdot2+1\). Only the sum of
the predecessor contributions decides the correction.

\section{Examples, cancellations, and transfinite endpoints}

Put \(\ep=\omega^{-1}\), and define the normal-form sum
\(u=\sum_{n=1}^{\infty}\ep^n\). The following values are direct
applications of \cref{thm:endpoint}. They also separate phenomena that
would be easy to miss in a monomial-only investigation.

\begin{center}
\renewcommand{\arraystretch}{1.28}
\begin{tabular}{@{}>{\raggedright\arraybackslash}p{0.29\textwidth}>{\raggedright\arraybackslash}p{0.20\textwidth}>{\raggedright\arraybackslash}p{0.45\textwidth}@{}}
\toprule
\textbf{Number} & \textbf{Birthday} & \textbf{Controlling feature}\\
\midrule
\(\ep^m\), \(1\leq m<\omega\) & \(\omega\cdot m\) & Dyadic monomial\\
\(3\ep^2\) & \(\omega\cdot2+2\) & Finite coefficient tail\\
\(\tfrac13\ep^2\) & \(\omega\cdot3\) & Non-dyadic last coefficient\\
\(1+\ep\) & \(\omega\) & Earlier finite length absorbed\\
\(\tfrac13+\ep\) & \(\omega+1\) & Successor endpoint correction\\
\(\tfrac13\ep+\ep^2\) & \(\omega\cdot2+1\) & Same correction above index 0\\
\(\omega^{-\omega}\) & \(\omega^2\) & Transfinite monomial index\\
\(u\) & \(\omega^2\) & Support has no maximum\\
\(u+\omega^{-\omega}\) & \(\omega^2+1\) & Cofinal support at last index\\
\(u+3\omega^{-\omega}\) & \(\omega^2+3\) & Cofinal endpoint with longer tail\\
\(u+\tfrac13\omega^{-\omega}\) & \(\omega^2+\omega\) & Non-dyadic terminal coefficient\\
\bottomrule
\end{tabular}
\end{center}

\subsection{Why a last-index formula without correction fails}
The tempting formula
\(\bd(x)=\omega\cdot a+d(c_a)-1\) would give birthday \(\omega\)
for \(\frac13+\ep\), instead of \(\omega+1\). It would also give
\(\omega^2\) for \(u+\omega^{-\omega}\), instead of \(\omega^2+1\).
These are counterexamples to that auxiliary formula, not counterexamples
to Gonshor's conjecture. The two failures have different causes, reflected
in the two clauses of \eqref{eq:correction}.

The second example is genuinely transfinite. No finite list of smaller
indices can be cofinal in \(\omega\). Testing arbitrarily many finite
truncations would therefore not, by itself, establish the birthday at the
new terminal index \(\omega\).

\subsection{Unbounded power-series support has a fixed birthday}
Every non-polynomial formal power series in \(\ep\), with real
coefficients, has an infinite support contained in \(\N\), hence
cofinal in \(\omega\). Therefore
\begin{equation}\label{eq:power-series}
 \bd\left(\sum_{n=0}^{\infty}c_n\ep^n\right)=\omega^2
 \quad\text{if infinitely many }c_n\ne0.
\end{equation}
This includes the coefficient sequences \(c_n=1\), \(c_n=1/n!\),
and any other real sequence with infinitely many nonzero entries.
No real-variable convergence hypothesis is involved.

An instructive cancellation uses
\[
 f=\sum_{n=0}^{\infty}\ep^n,\qquad
 g_-=1-\ep,\qquad g_+=1+\ep.
\]
Coefficientwise multiplication gives
\[
 fg_-=1,\qquad fg_+=1+2\sum_{n=1}^{\infty}\ep^n.
\]
The input birthdays are the same in the two products:
\(\bd(f)=\omega^2\) and \(\bd(g_-)=\bd(g_+)=\omega\).
The output birthdays are respectively 1 and \(\omega^2\).
Thus the product birthday is not determined by the two input birthdays,
even in this very small part of \(\A\).

\subsection{Cofinal supports at arbitrary limits}
Let \(\lambda\) be any nonzero limit ordinal and let \(S\subset\lambda\)
be cofinal, with no maximum. For arbitrary nonzero real coefficients,
\[
 \bd\left(\sum_{\alpha\in S}c_\alpha\omega^{-\alpha}\right)
    =\omega\cdot\lambda.
\]
Appending a last term at index \(\lambda\) with dyadic coefficient
\(c\) increases the birthday by \(d(c)\), not by \(d(c)-1\).
Appending a non-dyadic coefficient increases it by \(\omega\).

For \(\lambda=\omega^\rho\), the original birthday is
\(\omega^{1+\rho}\). When \(\rho=\omega\), this is
\(\omega^\omega\), not \(\omega^{\omega+1}\). This is one concrete
place where keeping ordinary left multiplication separate from natural
multiplication is indispensable.

\section{Set-sized local subrings and sharp birthday ceilings}

\subsection{A hierarchy of rings}
For an ordinal \(\rho>0\), set \(\kappa=\omega^\rho\) and define
\begin{equation}\label{eq:bounded-ring}
 \A_\kappa=\{x\in\A:S(x)\subseteq\kappa\}.
\end{equation}
Every such object is encoded by a function \(\kappa\to\R\), so this
is a set, unlike the full class \(\A\).

\begin{proposition}\label[proposition]{prop:bounded}
The set \(\A_\kappa\) is a subring of \(\A\), and
\begin{equation}\label{eq:ceiling}
 \max\{\bd(x):x\in\A_\kappa\}
      =\omega\cdot\kappa=\omega^{1+\rho}.
\end{equation}
A nonzero element attains this maximum exactly when its index support
is cofinal in \(\kappa\).
\end{proposition}
\begin{proof}
An ordinal below \(\omega^\rho\) has all Cantor exponents below
\(\rho\). Natural sum preserves this property, so the ring closure
proof applies without leaving \(\kappa\).

If a support has supremum \(\sigma<\kappa\), its birthday is at most
\(\omega\cdot(\sigma+1)<\omega\cdot\kappa\), since \(\kappa\)
is a limit ordinal and left multiplication by \(\omega\) is strictly
increasing in its right argument. A cofinal support has no maximum and
birthday \(\omega\cdot\kappa\) by \eqref{eq:no-maximum}. Such an
element exists: take coefficient 1 at every index below \(\kappa\).
Finally \(\omega\cdot\omega^\rho=\omega^{1+\rho}\).
\end{proof}

These ceilings are not obtained by restricting to numbers of birthday
\emph{strictly less} than a selected ordinal. They come from restrictions
on normal-form indices. The resulting ring can contain elements whose
birthday equals the stated ceiling, and products still remain inside.

\subsection{Units and a formal-power-series interpretation}
\begin{proposition}\label[proposition]{prop:local}
An element of \(\A_\kappa\) is a unit in that ring if and only if its
constant coefficient is nonzero. In particular \(\A_\kappa\) is a
local ring, with maximal ideal consisting of the elements having zero
constant coefficient.
\end{proposition}
\begin{proof}
The constant coefficient of a product is the product of the constant
coefficients: the only solution to \(\alpha\ns\beta=0\) is
\(\alpha=\beta=0\). Thus a zero constant coefficient prevents
invertibility.

Conversely write a series with nonzero constant coefficient as
\(c(1+v)\), where \(c\in\R\setminus\{0\}\) and \(v\) has zero
constant coefficient. Consider the formal geometric inverse
\[
 c^{-1}\sum_{n=0}^{\infty}(-v)^n.
\]
We verify local finiteness, rather than appealing to analytic convergence.
For \(\gamma=\sum_{j=1}^m\omega^{\rho_j}n_j\), put
\(\wt(\gamma)=\sum_j n_j\), and \(\wt(0)=0\). Natural sum adds
these weights. Every positive ordinal has weight at least 1. Consequently
a sum of \(n\) positive indices can equal \(\gamma\) only if
\(n\leq\wt(\gamma)\). For each such \(n\), distributing the finite
Cantor coefficient list among \(n\) slots gives only finitely many
ordered decompositions. Hence each coefficient of the geometric inverse
is a finite sum over finitely many values of \(n\).

All indices remain below \(\kappa\), by natural-sum closure.
Their union is a set of ordinals and hence well ordered, so the finite
occurrence just proved makes this a strong Hahn sum. To verify the
inverse identity at a fixed index \(\gamma\), take \(N=\wt(\gamma)\)
and use the finite polynomial identity
\[
 (1+v)\sum_{n=0}^{N}(-v)^n=1-(-v)^{N+1}.
\]
The coefficient at \(\gamma\) of every power with exponent greater
than \(N\) is zero. Consequently the infinite sum on the left has the
same \(\gamma\)-coefficient as this finite calculation, and the
remainder has zero coefficient there. This holds also at \(\gamma=0\),
where \(N=0\) and \(v\) has zero constant coefficient. Thus
\((1+v)\sum_{n\geq0}(-v)^n=1\), and the series is an inverse in
\(\A_\kappa\). The constant-coefficient map onto \(\R\) has as kernel
the stated ideal, and every element outside it is a unit. This proves
that the ideal is maximal and unique.
\end{proof}

To understand the algebra behind this proof, assign a formal variable
\(X_\xi\) to \(\omega^{-\omega^\xi}\), for each \(\xi<\rho\).
The Cantor normal form
\(\alpha=\sum_j\omega^{\xi_j}n_j\) corresponds to the ordinary monomial
\(\prod_j X_{\xi_j}^{n_j}\). Natural addition of indices becomes
multiplication of these monomials. Thus \(\A_\kappa\) is a
coefficientwise formal-series ring in these variables. When there are
infinitely many variables, this means arbitrary coefficient functions on
the finite-support monomials; it should not be confused without further
argument with a differently defined topological completion.

Every individual \(x\in\A\) belongs to some \(\A_\kappa\), by choosing
\(\rho\) with \(\sup S(x)<\omega^\rho\). This also avoids any need
to interpret the proofs as set-indexed sums over a proper class.

\section{Computation and reproducibility}\label{sec:computations}

\subsection{What the programs implement}
The companion module \texttt{code/ordinal\_series.py} represents ordinals
below \(\omega^\omega\) by finite tuples of nonnegative integers:
\[
 (a_0,a_1,\ldots,a_k)\longleftrightarrow
       \omega^k a_k+\cdots+\omega a_1+a_0.
\]
It distinguishes ordinary addition, natural addition, natural product,
ordinary left multiplication by \(\omega\), and the order type of an
ordinal interval. Coefficients are exact Python \texttt{Fraction} objects,
including both dyadic and non-dyadic rational numbers.

There are three routes to a finite-support product birthday.
The endpoint evaluator implements \cref{thm:endpoint}. The block evaluator
implements \eqref{eq:block-formula} directly, using shared-minus deletion
and the extra non-dyadic deletion. The four-coefficient evaluator implements
\cref{prop:four} without constructing the product. The verifier compares
all three, where applicable, against an explicitly constructed exact
coefficient convolution.

\subsection{Executed checks}
The default run used seed \texttt{20260920}. Its recorded result is
\texttt{PASS}; detailed output is in \path{verification.json}.
\begin{center}
\renewcommand{\arraystretch}{1.18}
\begin{tabular}{@{}lr@{}}
\toprule
\textbf{Check} & \textbf{Number executed}\\
\midrule
Dyadic birthday versus explicit finite sign string & 2,313\\
Ordinal interval concatenation identities & 55\\
Endpoint versus block birthday, exhaustive series family & 64,401\\
Exhaustive product pairs, unordered with repetition & 76,636\\
Additional seeded random product pairs & 20,000\\
Four-coefficient formula versus full convolution & 96,636\\
Named finite-support examples & 7\\
\bottomrule
\end{tabular}
\end{center}

The exhaustive birthday family uses support size at most three from ten
indices, including \(\omega\), \(\omega+1\), \(\omega\cdot2\), and
\(\omega^2\), with eight possible nonzero rational coefficients.
The exhaustive product family has 391 members: support size at most two
from five indices and six coefficient choices, including zero as the empty
series. Every unordered pair with repetition is checked. Random tests use
support size at most six and a seeded pool of 49 ordinal indices, with
rational coefficients whose denominators can contain odd factors.

For each product pair the program checks the main inequality, the equality
classification, the sharp support upper bound, the degree upper bound,
the strict gap when both birthdays are infinite, and agreement with the
four-coefficient formula. As an additional consistency check it verifies
the standard additive natural-sum bound on the same inputs.

\subsection{What these tests do not establish}
No infinite support is represented or executed. The code treats ordinal
indices with finite Cantor exponents, not arbitrary ordinals. The rational
coefficient tests do not enumerate arbitrary real coefficients. Most
importantly, agreement of the two birthday evaluators is not an
independent formal verification of the classical sign-expansion theorem:
both evaluators depend on that theorem, and both share the ordinal
arithmetic implementation.

The purpose of the tests is narrower and useful: they expose off-by-one
errors, ordinary/natural operation confusion, failure to handle
non-dyadic coefficients, missed coefficient cancellation, and mistakes in
the proposed equality condition. The transfinite proof in
\cref{sec:endpoint}, including its cofinal endpoint case, is the reason
the mathematical statement covers infinite supports.

\subsection{Reproducing the run}
Only Python 3.10 or later and its standard library are needed:
\begin{lstlisting}
python code/verify.py
python code/verify.py --seed 20260920 --random-pairs 20000
\end{lstlisting}
The default command writes \path{verification.json}. To preserve the
shipped record, use a separate output file:
\begin{lstlisting}
python code/verify.py --output verification-rerun.json
\end{lstlisting}
A different seed or a larger random sample can be selected explicitly.
Failure raises an assertion with the offending pair or series. The code
is research verification code, not a general-purpose surreal package.

\section{Assessment of the attempt and the remaining general problem}

\subsection{What has actually been proved}
The endpoint theorem gives exact birthdays on a ring of real-coefficient
series with arbitrary set-sized ordinal-index support. It is not merely a
calculation for powers \(\omega^{-n}\), a finite-support observation, or
an estimate for individual monomials. The product theorem covers the
whole ring, including cancellation and transfinite accumulation of support.
Its strictness conclusion and its equality classification are consequences
of a leading-Cantor-degree gap, not of numerical experimentation.

The associated four-coefficient formula makes a finite-support output
birthday substantially easier to compute than the output normal form.
The local-ring interpretation and the attained birthday ceilings show
that the domain has familiar and nontrivial algebraic structure.

The main claims are deductions made in this investigation from classical
sign-expansion theory. The literature search did not establish whether
the exact endpoint formula or these specific deductions have appeared
elsewhere. They should therefore be described as the results of this
attempt, not as certified first discoveries.

\subsection{Why the proof does not extend automatically}
For a general negative surreal exponent, the sign string can contain plus
signs. For example, \(-\tfrac12\) has sign string \(-+\), whereas
\(-\alpha\) for an ordinal \(\alpha\) has only minuses. In the former
setting the exponent blocks are no longer uniformly \(\omega\)-long.
Reduced sign strings can have several retained regions rather than one
terminal interval. The proof of \eqref{eq:block-formula} in its simplified
form, and therefore the endpoint induction, no longer applies.

Allowing positive exponents changes the picture even more visibly.
The numbers \(x=y=\omega\) satisfy
\[
 \bd(xy)=\omega^2=\bd(x)\np\bd(y),
\]
although both birthdays are infinite. Thus the strictness theorem and
equality classification proved for \(\A\) cannot be generalized to all
of \(\No\) without a different statement. This example does not refute
the general conjecture; it isolates the limitation of the stronger
special-case result.

\subsection{A concrete next extension}
A natural next domain would allow exponents whose sign strings have one
controlled plus block after an initial minus block. A successful extension
would need a replacement for the single-interval calculation: it would
track which exponent signs survive reduction, the ordinal sizes of their
expanded blocks, and how these data behave under addition of exponents.
The four-coefficient shortcut suggests separating the analysis of support
endpoints from the analysis of coefficient tails, but it does not solve
that new block-comparison problem.

The general conjecture remains a precise further target. The result of
this attempt is a complete special-case proof with explicit boundaries,
not a proof obtained by silently identifying all negative exponents with
negative ordinals or by treating finite tests as transfinite evidence.

\appendix
\section{A compact proof audit}

The following checks are useful when reading or extending the argument.

\paragraph{Canonical objects.}
The quantity \(\bd(x)\) is the length of the canonical sign expansion.
Recursive form generation is never substituted for it.

\paragraph{Ordinal operations.}
Normal-form indices convolve by \(\ns\). The conjectured upper bound
uses \(\np\). Concatenated sign lengths use ordinary \(+\) and
ordinary left multiplication \(\omega\cdot\alpha\). In
\(\omega^{1+\rho}\), the exponent addition is ordinary.

\paragraph{Coefficient types.}
Dyadic status is checked on the actual nonzero coefficient after all
contributions at an index have been added. Products of non-dyadic reals
can be dyadic, and a dyadic scalar can turn a non-dyadic real into a
dyadic one. Both possibilities are included.

\paragraph{Limit stages.}
The limit argument takes the supremum of lengths of concatenated sign
prefixes. It does not assume any general continuity of surreal addition
or multiplication. A last index cofinal with its predecessors contributes
an entire coefficient string, explaining the correction by one.

\paragraph{Set-sizedness and local finiteness.}
Every support is a set. The natural-sum fiber count
\eqref{eq:fiber-count} makes each convolution coefficient finite.
For the geometric inverse, the finite Cantor weight bounds the number
of factors that can contribute at a prescribed index.

\paragraph{Finite endpoint cancellation.}
With finite nonempty supports, the maximum product index has a unique
contributing pair and cannot cancel. At its predecessor there can be two
contributions; their actual sum determines dyadic status. The
four-coefficient formula is a mathematical dependence statement for real
inputs and an executable algorithm for the exact rational inputs in the code.

\paragraph{Equality.}
The cases ``both finite'', ``one finite'', and ``both infinite'' are
exhaustive because finite-birthday surreals are precisely dyadic reals.
Strictness in the last case follows from a full leading-exponent gap;
strictness in the mixed case follows from the leading coefficient.

\section{Source and artifact notes}

The historical conjecture is documented by \cite{OPG} and the later
first-person account \cite{EhrlichMO}. The latter is the source used here
for the reported scope of the van den Dries--Ehrlich monomial result.
The original 2001 article was not used as a directly inspected proof
source in preparing this draft.

The actual sign-expansion ingredient was checked in
\cite[Definition~2.20 and Theorem~2.21, printed p.~13]{BournezGuilmant}.
Gonshor's book is credited as the original source of that theorem; the
version used in this derivation is the explicitly identified restatement.
The normal-form background is as in \cite{Berarducci}. The distinction
between forms and values, and the finite dyadic birthday formula, are
also visible in \cite{Roughan}.

The archive contains this \LaTeX\ source, its compiled PDF, the
\texttt{references.bib} bibliography, two Python files,
\path{verification.json}, a README, and build instructions. It does not
redistribute the cited books or papers. The mathematical article is the
primary result; the code provides reproducible finite checks with the
limitations stated in \cref{sec:computations}.

\begin{thebibliography}{Gon86}
\raggedright
\bibitem[Ber20]{Berarducci}
Alessandro Berarducci.
\emph{Surreal numbers, exponentiation and derivations}.
2020. arXiv:2008.06878v1, especially Sections 5, 6, and 8--10.
\url{https://arxiv.org/abs/2008.06878}.

\bibitem[BG22]{BournezGuilmant}
Olivier Bournez and Quentin Guilmant.
\emph{Surreal fields stable under exponential and logarithmic functions}.
2022. arXiv:2201.08199v1, especially Definition 2.20 and Theorem 2.21,
printed p.~13. \url{https://arxiv.org/abs/2201.08199}.

\bibitem[Ehr24]{EhrlichMO}
Philip Ehrlich.
Answer to ``Birthday of combinatorial game product''.
\emph{MathOverflow}, question 476824, 13 August 2024, edited the same day.
Consulted 20 September 2026.
\url{https://mathoverflow.net/questions/476824}.

\bibitem[Gon86]{Gonshor}
Harry Gonshor.
\emph{An Introduction to the Theory of Surreal Numbers}.
London Mathematical Society Lecture Note Series, vol.~110.
Cambridge University Press, 1986.
Original source credited for the conjecture and the normal-form
sign-expansion theorems; the latter are used via [BG22].

\bibitem[L\'an12]{OPG}
Luk\'a\v{s} L\'ansk\'y.
``Length of surreal product''.
\emph{Open Problem Garden}, 7 April 2012. Consulted 20 September 2026.
\url{https://www.openproblemgarden.org/op/length_of_surreal_product}.

\bibitem[Rou18]{Roughan}
Matthew Roughan.
\emph{Surreal Birthdays and Their Arithmetic}.
2018. arXiv:1810.10373v2.
\url{https://arxiv.org/abs/1810.10373}.
\end{thebibliography}
\end{document}
